\documentclass[a4paper]{report}

\usepackage{graphicx}
\usepackage{float}
\usepackage{parskip}
\usepackage{geometry}
\usepackage[linesnumbered, ruled]{algorithm2e}
\usepackage{array}
\usepackage{listings}
\usepackage{xcolor}
\usepackage{amsmath}
\usepackage{caption}
\usepackage{varioref}
\usepackage[hidelinks]{hyperref}
\usepackage{cleveref}

% Preamble

% Numbering on subsubsections
\setcounter{secnumdepth}{3}

% Prevent tabular p{} spacing out text
\newcolumntype{P}[1]{>{\RaggedRight\arraybackslash}p{#1}}

% Comments for pseudocode
\SetKwComment{tcp}{// }{}
\newcommand\commentfont[1]{\textcolor{blue}{#1}}
\SetCommentSty{commentfont}

% Use better font encoding
\usepackage[T1]{fontenc}

% Shrink the huge default margins around the sides of the document
\newgeometry{vmargin=1in, hmargin=1in}

% Command to restore margins
\newcommand\restoremargins{\pagebreak\restoregeometry}

% Chapter author command
\makeatletter
\newcommand{\chapterauthor}[1]{%
	{\parindent0pt\vspace*{-25pt}%
		\linespread{1.1}\large\scshape#1%
		\par\nobreak\vspace*{25pt}}
	\@afterheading%
}

% Section author command
\newcommand{\sectionauthor}[1]{%
	{\parindent0pt\vspace*{0pt}%
		\linespread{1.1}\large\scshape#1%
		\par\nobreak\vspace*{10pt}}
	\@afterheading%
}

% Extend footnote rule
\renewcommand\footnoterule{%
	\kern-3\p@
	\hrule\@width \textwidth
	\kern2.6\p@}
\makeatother
\makeatother

% Colour definitions
\definecolor{codegreen}{rgb}{0,0.6,0}
\definecolor{codegrey}{rgb}{0.5,0.5,0.5}
\definecolor{codepurple}{rgb}{0.58,0,0.82}
\definecolor{backcolour}{rgb}{0.95,0.95,0.92}

% Code style definition
\lstdefinestyle{code}{
	language=java,
	tabsize=4,
	commentstyle=\color{codegreen},
	backgroundcolor=\color{backcolour},
	keywordstyle=\color{magenta},
	numberstyle=\tiny\color{codegrey},
	stringstyle=\color{codepurple},
	basicstyle=\ttfamily, 
	breakatwhitespace=false,
	breaklines=true,
	numbers=left,
	numbersep=5pt,
	gobble=16,
	showstringspaces=false,
	captionpos=t
}
\lstset{style=code}

% Code caption formatting
\renewcommand\lstlistingname{}
\DeclareCaptionFormat{listing}{\rule{\dimexpr\textwidth\relax}{0.4pt}\vskip1pt#1#2#3}
\captionsetup[lstlisting]{format=listing,singlelinecheck=false, margin=0pt, font={sf}}

% Make command for blue hrefs
\newcommand{\link}[2]{\color{blue}\href{#1}{#2}\color{black}}

% Title and author are used by \maketitle on the title page
\title{Beginner's Guide To Java Unit Tests - Spike Work}
\author{Group 6}

\begin{document}
	\maketitle
	\tableofcontents
	\newpage
	
	\chapter{Introduction}\label{chap:Java Unit Tests}
	\chapterauthor{Niall Perry}

JUnit is a widely used framework for unit testing in Java. In this guide, I will show how to perform general Java unit tests and then Java unit tests for JavaFx.

\section{General J unit tests}\label{sec:introduction}
        What is Unit Testing?
A procedure to validate individual units of the source code

Example: A procedure, method, or class

Validating each individual piece reduces errors when integrating the pieces together later.

A recommendation is to have tests that focus on
boundaries of allowed data and where erroneous data are handled correctly...

If a method takes data 0.. 10 What should you test?
Answer: -1, 0, 10, 11


	\section{Outline of a Test class and Test methods}\label{chap:Test Class}
	\title{Test class}
\begin{lstlisting}
    
                        package group.tests;
                    import org.junit.jupiter.api.Assertions;
                    import org.junit.jupiter.api.BeforeEach;
                    import org.junit.jupiter.api.Test;
                    import group6.*;
                    
                    public class AlbumTest {
                    private Album album;
                    
                    @BeforeEach
                    public void setup(){
                    album = new Album();
                    }
                    
                    @Test
                    public void testAddOneSongAndFind() {
                    // Test code here
                    }
                    }
\end{lstlisting}
Tests method marked with @Test is run…

@BeforeEach methods are always run before each
@Test method…

There are also @AfterEach, @BeforeAll, @AfterAll

\label{chap:Test methods}
\title{Test methods}
\begin{lstlisting}
                        @Test
                    public void testAddOneSongAndFind() { // [1]
                    
                    // Set up test fixture
                    Contact contact = new Contact("Hey Jude", "The Beatles"); // [2]
                   
                    // Exercise the program under test [3]
                    album.add(contact);
                    Song foundSong = album.search("Hey Jude");
                   
                    // Check expected data against data returned [4]
                    Assertions.assertEquals(song, foundSong,
                    "Incorrect song found!");
                    }
\end{lstlisting}
Things to note:

[1] Meaningful name

[2] Create the test data: called the test fixture

[3] Do something interesting and get a result

[4] Did we get the right result?

\textbf{Mantra: Arrange, Act, Assert}

\chapter{How to make Tests using Inteli-J} \label{Step 1}
We will create a JUnit test to check that songs can be added to an album.

\section{Step 1}\label{sec:Step 1}
\textbf{Create package}

1. Right-click on src

2. Select New

3. Select Package


\begin{figure}[H]
                        \begin{center}
                        \includegraphics[width=0.67]{Screenshot 2025-03-08 132341.png}
                        \end{center}
                        \caption{a visual guide}
                    \end{figure}


\section{Step 2}\label{sec:Step 2}
\textbf{Create test class}

1. Right click on the tests package

2. Select New

3. Select Class

\begin{figure}[H]
    \begin{center}
    \includegraphics[width=0.80]{Screenshot 2025-03-08 134743.png}
    \end{center}
    \caption{visual guide}
\end{figure}
\newpage
\section{Step 3}\label{sec:Step 3}
\textbf{Adding to test class}

1. Write some code

2. Import the JUnit library

\begin{figure}[H]
    \begin{center}
    \includegraphics[width=0.8\textwidth]{Screenshot 2025-03-08 144856.png}
    \end{center}
    \caption{visual guide}
\end{figure}

\section{Step 4}\label{sec:Step 4}
\textbf{Run the test class}

1. run code

\begin{figure}[H]
    \begin{center}
file:///C:/Users/pezuk/OneDrive - Aberystwyth University/Year 2/Software engineering/Practice Junit/Junit_tests.aux
file:///C:/Users/pezuk/OneDrive - Aberystwyth University/Year 2/Software engineering/Practice Junit/Junit_tests.dvi
file:///C:/Users/pezuk/OneDrive - Aberystwyth University/Year 2/Software engineering/Practice Junit/Junit_tests.log
file:///C:/Users/pezuk/OneDrive - Aberystwyth University/Year 2/Software engineering/Practice Junit/Junit_tests.out
file:///C:/Users/pezuk/OneDrive - Aberystwyth University/Year 2/Software engineering/Practice Junit/Junit_tests.synctex.gz
file:///C:/Users/pezuk/OneDrive - Aberystwyth University/Year 2/Software engineering/Practice Junit/Junit_tests.tex
file:///C:/Users/pezuk/OneDrive - Aberystwyth University/Year 2/Software engineering/Practice Junit/Junit_tests.toc
    \includegraphics[width=0.67]{Screenshot 2025-03-08 155244.png}
    \end{center}
    \caption{visual guide}
\end{figure}

Here is the expected output...

\begin{figure}[H]
    \begin{center}
    \includegraphics[width=0.50]{Screenshot 2025-03-11 002045.png}
    \end{center}
    \caption{Result of action}
\end{figure}
\newpage
\section{Step 5}\label{sec:Step 5}
\textbf{Add proper test code }

For example:


\begin{lstlisting}
                                 public class Album_test {
                    
                        public Album album;
                    // intialises test
                        @BeforeEach
                        void setUp() {
                        album = new Album("Michael Jackson", 1982, "Thriller");}
                    
                    // example of testing contructor
                        @Test
                        void testConstructor() {
                            assertEquals("Thriller", album.getAlbumTitle());
                            assertEquals("Michael Jackson", album.getArtist());
                            assertEquals(1982, album.getYear());
                            assertTrue(album.getTracks().isEmpty(), "Track list should be empty initially");}
                    
                    // example of testing track addition
                        @Test
                        void testAddTrack() {
                            album.addTracks("Billie Jean");
                            album.addTracks("Beat It");
                            List<String> tracks = album.getTracks();
                            assertEquals(2, tracks.size());
                            assertTrue(tracks.contains("Billie Jean"));
                            assertTrue(tracks.contains("Beat It"));
                        }
                    }
    
\end{lstlisting}
*** This can only be done if the method tested exists.

\section{Step 6}\label{sec:Step 6}
\textbf{finally}

1. compile the code

\begin{figure}[H]
    \begin{center}
    \includegraphics[width=0.5]{Screenshot 2025-03-11 031534.png}
    \end{center}
    \caption{the result}
\end{figure}

\chapter{Java FX - J Unit tests}\label{sec: Java FX J unit tests}

\section{Best practices} \label{sec: Best practices}
\textbf{1. Use TestFX for UI Interactions}

\begin{lstlisting}
                import org.testfx.api.FxRobot;
                
                class MainAppTest extends ApplicationTest {
                
                    private Button button;
                
                    @Override
                    public void start(Stage stage) {
                        button = new Button("Click Me");
                        button.setOnAction(event -> button.setText("Clicked!"));
                
                        stage.setScene(new javafx.scene.Scene(new javafx.scene.layout.StackPane(button), 300, 200));
                        stage.show();
                    }
                
                    @Test
                    void testButtonClick(FxRobot robot) {
                        assertEquals("Click Me", button.getText()); // Verify initial state
                
                        robot.clickOn(button); // Simulate a click
                
                        assertEquals("Clicked!", button.getText()); // Verify updated text
                    }
                }
\end{lstlisting}

JavaFX involves UI components that require a JavaFX thread to function correctly. TestFX provides utilities to interact with JavaFX elements.

For example, using FxRobot
\newpage

\textbf{2. Always Run JavaFX Tests on the JavaFX Thread}
\begin{lstlisting}
                    import javafx.application.Platform;
                import org.junit.jupiter.api.Test;
                import static org.junit.jupiter.api.Assertions.*;
                
                class JavaFXThreadTest {
                
                    @Test
                    void testFXThreadExecution() {
                        Platform.runLater(() -> {
                            System.out.println("Running on JavaFX Thread!");
                            assertTrue(Platform.isFxApplicationThread());
                        });
                    }
                }
\end{lstlisting}

\textbf{3. Use @Start for Stage Initialization}

Instead of setting up the UI in @BeforeEach, use @Start in TestFX.

@Start automatically launches a JavaFX stage, avoiding manual setup.

The start(Stage stage) method ensures UI elements are created properly.

\begin{lstlisting}
                 public void start(Stage stage) {
                        button = new Button("Click Me");
                        button.setOnAction(event -> button.setText("Clicked!"));
                
                        stage.setScene(new javafx.scene.Scene(new javafx.scene.layout.StackPane(button), 300, 200));
                        stage.show();
\end{lstlisting}

\textbf{4. Use Timeouts for Asynchronous UI Operations}

Some JavaFX operations run asynchronously (e.g., animations, background tasks). To avoid flaky tests, use TestFX’s wait conditions.

\begin{lstlisting}
                import org.testfx.util.WaitForAsyncUtils;
                import java.util.concurrent.TimeUnit;
            
                WaitForAsyncUtils.sleep(1, TimeUnit.SECONDS); // Wait before checking
                assertEquals("Clicked!", button.getText());
\end{lstlisting}

\textbf{5. Use lookup() for Selecting UI Elements}

\begin{lstlisting}
                    @Test
                    void testButtonClick(FxRobot robot) {
                    robot.clickOn(".button"); // Find button by CSS class
                    assertThat(lookup(".button").queryButton()).hasText("Clicked!");
                }
\end{lstlisting}

\textbf{6. Use Assertions for UI Elements}

Use TestFX Assertions to validate UI state.

\begin{lstlisting}
                import static org.testfx.assertions.api.Assertions.assertThat;
            
                assertThat(button).hasText("Clicked!");
                assertThat(label).isVisible();
                assertThat(textField).hasText("Hello World");
\end{lstlisting}    

\end{document}